% =====================================================================
% APPENDIX
% =====================================================================

\section{Full training details}
\label{app:training}

\subsection{Shared CNN setup (SGD and SGLD comparison)}

\paragraph{Architecture.}
A ``$k$-width'' CNN: 4 convolutional blocks (channels $3 \to k \to 2k \to 4k
\to 8k$), each block consisting of a $3\times 3$ convolution (padding 1),
group normalization (one group), ReLU, and max pool with strides $[1, 2, 2, 8]$, so a
$32\times 32$ input collapses to $1\times 1$ after the four pools. A final
linear head maps $8k \to 10$. Models are trained from scratch with random
initialization; no pretraining and no data augmentation.

\paragraph{Dataset.}
CIFAR-10 with pixels rescaled to $[0,1]$ and standardized per-channel
using train-set mean and standard deviation. The first $n_{\mathrm{train}}
= 50{,}000$ images form the training set; the standard $10{,}000$-image
test images are unchanged. Labels are corrupted by replacing each label
with a uniformly random class with probability equal to the target noise
level, using a fixed RNG (\texttt{np.random.RandomState(seed=0)}) so that
the corruption pattern is identical across seeds and widths.

\paragraph{Optimization.}
SGD with no momentum and no weight decay; cross-entropy loss on the noisy
labels; fp32 throughout with TF32 disabled. Learning rate follows the
schedule
\[
\mathrm{lr}(t) \;=\; \mathrm{lr}_0 \,/\, (1 + t/T_{\mathrm{decay}})^{0.55},
\]
with $\mathrm{lr}_0 = 0.1$ and $T_{\mathrm{decay}} = 512$.
Minibatches of 128 examples are sampled uniformly with replacement.

\paragraph{Gaussian-prior SGLD.}
Every trainable parameter has an independent Gaussian prior. Convolutional
weights have mean zero and variance $2/\mathrm{fan\_in}$; linear weights
have mean zero and variance $1/\mathrm{fan\_in}$; biases have mean zero and
variance $0.25$; group-normalization scales have mean one and variance
$0.25$. For mean minibatch cross-entropy $\ell_B$, the update is
\[
\theta_{t+1}=\theta_t-\mathrm{lr}(t)
\left[\nabla\ell_B(\theta_t)-\frac{1}{n}\nabla\log\pi(\theta_t)\right]
+\sqrt{\frac{2\,\mathrm{lr}(t)}{n}}\,\xi_t,\qquad \xi_t\sim N(0,I).
\]
This is the temperature-one SGLD update \citep{welling2011sgld}, with
posterior target proportional to the likelihood times the stated prior.
Group normalization makes each prediction independent of the other
examples in its minibatch.

\paragraph{Per-seed protocol.}
Each seed receives its own \texttt{torch.manual\_seed(seed)} for weight
initialization and a CUDA \texttt{Generator(seed+123456)} for batch
sampling. Inputs and (corrupted) labels are identical across seeds; only
model initialization, minibatch order, and SGLD noise vary. For each (width $k$, seed)
pair we save the post-training softmax probabilities on the test set,
along with train cross-entropy and 0/1 error against the noisy training
labels, test cross-entropy against \emph{noisy} test labels (corrupted
with the same probability as the training labels), and test 0/1 error
against the \emph{clean} test labels. Test corruption uses
\texttt{np.random.RandomState(seed=12345)}, independently of training
corruption. SGLD discards the first $100{,}000$ steps and retains
predictive probabilities every $10{,}000$ steps thereafter, giving 10
draws per chain. Within-chain averages weight each retained draw by the
sum of step sizes since the preceding retained draw (or burn-in endpoint);
the final ensemble weights the eight chains equally. SGD predictions use
the final iterate. Log-loss evaluation floors assigned probabilities at
$10^{-12}$; the largest clipped fraction across widths is $0.504\%$ for
20\% noise and $0.223\%$ for 100\% noise, in the individual SGD predictions.

\paragraph{Test labels: clean vs.\ noisy.}
We use noisy test labels for cross-entropy and clean test labels for 0/1
error. For 0/1 error, expected noisy loss is an affine function of the
clean loss: writing the noise level as $\eta\in[0,1]$ and the number of
classes as $K=10$, for any classifier $\hat Y(X)$
\begin{equation}
\mathbb{E}_{\widetilde{Y}\mid Y}\!\left[\mathbb{1}\{\hat Y(X) \ne \widetilde{Y}\}\right]
\;=\; (1-\eta)\,\mathbb{1}\{\hat Y(X) \ne Y\}
\;+\; \eta\cdot\frac{K-1}{K},
\end{equation}
which is strictly-increasing affine in the clean 0/1 loss whenever
$\eta<1$, so the curve shape (peak and descent) is identical under either
choice up to an affine transformation. We use the clean 0/1 error for
interpretability since 0/1 is reported for reference rather than as a core
theoretical metric. For cross-entropy, by contrast, no such relation
holds:
\begin{equation}
\mathbb{E}_{\widetilde{Y}\mid Y}\!\left[-\log q(\widetilde{Y}\mid X)\right]
\;=\; (1-\eta)\,\bigl(-\log q(Y\mid X)\bigr)
\;-\; \frac{\eta}{K}\sum_{c=1}^{K} \log q(c\mid X),
\end{equation}
whose second term depends on the full predictive distribution rather
than only on $q(Y\mid X)$, so noisy cross-entropy is not a monotonic
function of clean cross-entropy. We therefore evaluate cross-entropy on
the noisy test labels, which are the relevant labels for the data
distribution $P^*$ in our decomposition.

\subsection{CIFAR-10 experiments}

\paragraph{20\% noise.}
Used in Figures~\ref{fig:ensemble-collapse} (Right), \ref{fig:ensemble}
(top row).
$n_{\mathrm{train}} = 50{,}000$, label noise $= 20\%$;
widths $k \in \{1, 2, 3, 4, 5, 6, 8, 10, 14, 20, 32, 48, 64\}$;
batch size $128$; $200{,}000$ gradient steps;
$K=8$ SGD seeds and $K=8$ SGLD chains per width;
\texttt{torch.compile} enabled.

\paragraph{100\% noise.}
Used in Figure~\ref{fig:ensemble} (bottom row).
$n_{\mathrm{train}} = 50{,}000$, label noise $= 100\%$ (all labels random);
widths $k \in \{1, 2, 3, 4, 5, 6, 8, 10, 14, 20, 32, 48, 64, 96, 128\}$.
All other settings as in the 20\% noise experiment.

\paragraph{SGD ensembling.}
Figure~\ref{fig:gap-vs-seeds} uses archived SGD predictions from 64 seeds
per width under 20\% noise and 16 under 100\% noise. The architecture is
as above with batch normalization in place of group normalization;
training and test inputs are standardized separately. SGD has no momentum
or weight decay and runs for $200{,}000$ steps with batch size 128 and
$\mathrm{lr}(t)=0.1/\sqrt{1+\lfloor t/512\rfloor}$.
Minibatches follow shuffled passes through the training set, with each
seed used for both initialization and the CUDA shuffling generator.
Widths are as above, with $\{192,256,384\}$ additionally included under
100\% noise. Label corruption and test metrics are as above. Ensembles
average the final-iterate class probabilities across seeds.

\subsection{Random forests on MNIST}

\paragraph{Dataset.}
MNIST with pixels divided by $255$ (range $[0,1]$); the official
$60{,}000$-train / $10{,}000$-test split is used. We uniformly subsample
$n = 10{,}000$ training images without replacement (RNG seed $42$) and
evaluate on the full $10{,}000$-image test set. Targets are one-hot
encoded over $K = 10$ classes.

\paragraph{Model.}
We use \texttt{sklearn.ensemble.RandomForestRegressor} with
\texttt{bootstrap=False} and \texttt{max\_features} $= \lfloor\sqrt{784}\rfloor = 28$.
The train and test curves use the archived single-forest results.

\paragraph{Two-phase complexity sweep.}
Along the $x$-axis $N^{\mathrm{max}}_{\mathrm{leaf}} / N_{\mathrm{tree}}$:
Phase~1 fixes $N_{\mathrm{tree}} = 1$ and varies
$N^{\mathrm{max}}_{\mathrm{leaf}}$ from small values up to $n=10^4$
(passing through the interpolation threshold); Phase~2 fixes
$N^{\mathrm{max}}_{\mathrm{leaf}}$ at the interpolation threshold and
varies $N_{\mathrm{tree}} \in \{2, 3, 5, 7, 10, 12, 15, 20\}$. We run two
versions: one with the clean MNIST labels (Phase~2 at
$N^{\mathrm{max}}_{\mathrm{leaf}}=2000$, just past the empirical
interpolation point), and one with $100\%$ random labels (Phase~2 at
$N^{\mathrm{max}}_{\mathrm{leaf}}=n=10{,}000$, since fully fitting random
labels requires one leaf per training point).

\paragraph{BART comparison.}
The comparison uses a sum of $m$ regression trees, each capped at $L$
leaves, with one-hot Gaussian regression variance $\sigma^2=0.1$ and
independent Gaussian leaf values of variance $0.0625/m$ and mean zero.
The split probability is $0.5$ at every depth, with uniform feature and
available-threshold proposals. Each of 16 chains uses 3,000 burn-in
sweeps and retains 25 draws, one every 120 sweeps; the number of structure
moves per tree per sweep is
$\max(1,\mathrm{round}((30L+3000)/3000))$.
The BART curve averages regression outputs across draws and chains, then
clips and normalizes the ten scores. This transformation is not the
categorical posterior predictive of the Gaussian regression model.

\paragraph{Metrics.}
We report squared loss (used only to identify the interpolation
threshold as the smallest $x$-value at which the average train squared
loss first falls below $10^{-4}$), zero-one loss, and cross-entropy.
Cross-entropy is computed on the row-renormalized, $\varepsilon$-clipped
one-hot-mean prediction with $\varepsilon = 10^{-12}$, ensuring a proper
probability vector.

\section{Compute resources}
\label{app:compute}

The CNN reruns used four AWS Spot instances, each with eight NVIDIA H100
80\,GB GPUs and a 200\,GB root volume. The 448 main runs and 48 sampling
checks used approximately 145 allocated H100 GPU-hours, including
provisioning and shutdown. Recorded training times were 6.9--62.6 minutes
per main run and 14.5--121.6 minutes per sampling check.
The archived forest and BART experiments ran on CPU; their hardware and
runtime details were not reconstructed. Hardware and runtime details for
the archived SGD runs in Figure~\ref{fig:gap-vs-seeds} were also not
reconstructed. These figures also exclude earlier exploratory runs.

\section{Notation and empirical quantities}
\label{app:interpretation}

\begin{center}
\begin{tabular}{lp{0.68\linewidth}}
\toprule
Quantity & Interpretation and measurement \\
\midrule
$L_n(P^*)$ & Expected plug-in test log-loss; estimated by the mean SGD test loss. \\
$Q^\pi_{n+1}$ & Bayes posterior predictive under the stated prior; approximated by SGLD. \\
(I) & Change in plug-in risk from distribution shift; zero under uniform random labels. \\
(II), (II') & Minimax risk and prior cost; not estimated in the CNN experiments. \\
(III) & Worst-case MLE-vs-Bayes risk difference; not directly estimated. \\
Empirical gap & SGD test loss minus SGLD ensemble test loss; the expectation lower bound requires unbiased posterior prediction. \\
\bottomrule
\end{tabular}
\end{center}

\section{Sampling checks and ensemble control}
\label{app:sampling}

We run four seeds at widths $\{8,14,64\}$ under 20\% noise and
$\{10,20,128\}$ under 100\% noise. The longer-run check uses 400,000
steps, discards 200,000, and retains a draw every 20,000 steps. The
smaller-step check uses 400,000 steps with
$\mathrm{lr}(t)=0.05/(1+t/1024)^{0.55}$, discards 200,000 steps, and
retains a draw every 20,000 steps, matching the main run's effective
integration time. Both comparisons use the same four seeds as their
200,000-step reference.

The largest absolute change in ensemble log-loss is at most $0.0186$ nats under
20\% noise and $0.000319$ nats under 100\% noise. Mean predictive total
variation from the reference ranges from $0.065$ to $0.117$ and from
$0.008$ to $0.018$, respectively. The qualitative loss comparison is
preserved, but predictive differences and retained-window drift remain;
these checks do not establish convergence.

\begin{figure}[t]
  \centering
  \includegraphics[width=\linewidth]{figures/sampling_checks.pdf}
  \caption{CIFAR-10 with 20\% label noise (top) and 100\% label noise
  (bottom). Each ensemble uses the same four seeds and 10 retained draws
  per chain. \textbf{Left:} Ensemble test cross-entropy for the main run,
  longer run, and smaller-step run. \textbf{Right:} Mean predictive total
  variation from the main-run ensemble at the same width.}
  \label{fig:sampling-checks}
\end{figure}

\begin{figure}[t]
  \centering
  \includegraphics[width=\linewidth]{figures/ensemble_comparison.pdf}
  \caption{CIFAR-10 with 100\% label noise (top) and 20\% label noise
  (bottom), with $K=8$ seeds or chains per width.
  \textbf{Left:} Per-seed SGD test cross-entropy (orange; mean over seeds,
  band $\pm 1$ s.d.), the probability-averaged SGD ensemble (blue, dotted),
  and the SGLD ensemble (green, dashed).
  \textbf{Right:} The per-seed SGD mean minus each ensemble's test
  cross-entropy at the same width.}
  \label{fig:ensemble-control}
\end{figure}
